%!TEX TS-program = lualatex
%!TEX encoding = UTF-8 Unicode

\documentclass[11pt, addpoints]{exam}

%\printanswers


\usepackage[left=1in,right=1.25in,top=1in]{geometry}     

\usepackage{fontspec}
\def\mainfont{Linux Libertine O}
\setmainfont[Ligatures={TeX, Common}, BoldFont={* Bold}, ItalicFont={* Italic}, Numbers={Proportional, OldStyle}]{\mainfont}
\setsansfont[Scale=MatchLowercase]{Linux Biolinum O} 
\usepackage{microtype}

\usepackage{graphicx}
	\graphicspath{%
	{/Users/goby/Pictures/teach/348/exams/}}%

% Get the last two digits of the year for the header below.
\def\short#1{\csname @gobbletwo\expandafter\endcsname\number#1}

\pagestyle{headandfoot}
\firstpageheader{Marine Biology, Exam 3, \textsc{s}\short{\year}.}{}{%
	\ifprintanswers\textbf{KEY}\else Name: \enspace \makebox[2.5in]{\hrulefill}\fi}
\runningheader{}{}{\small{pg. \thepage}}
%\runningheadrule

\footer{\iflastpage{\hspace{-1.2cm}\small \rule{0.6in}{0.4pt} / \numpoints\ points}{}}{}{}%{\makebox[0.75in]{\hrulefill}}
\extrafootheight{-0.5in}

\pointsinmargin
\marginpointname{ pts}
\marginbonuspointname{ pts \textsc{e.c.}}

\usepackage{multicol}
\raggedcolumns

%% Remove comment %% to print answer key.
\unframedsolutions
\renewcommand{\solutiontitle}{}
\SolutionEmphasis{\bfseries}


\newcommand*\AnswerBox[2]{%
    \parbox[t][#1]{0.92\textwidth}{%
    \begin{solution}#2\end{solution}}
    \vspace{\stretch{1}}
}

\newenvironment{AnswerPage}[1]
    {\begin{minipage}[t][#1]{0.92\textwidth}%
    \begin{solution}}
    {\end{solution}\end{minipage}
    \vspace{\stretch{1}}}

\newlength{\basespace}
\setlength{\basespace}{5\baselineskip}


\begin{document}

\begin{questions}

\question[3]\label{rocky_intertidal_zones}
Name the three zones of the rocky intertidal, from mean low tide to mean high tide.

\AnswerBox{0.1\textheight}{infralittral zone, midlittoral zone, supralittoral fringe.}


\question[4]
\begin{multicols}{2}
	What is the name of the tidal pattern shown at right. Tell how you know.
	
	\begin{solution}
		Mixed semidurnal. Has two high tides of unequal height in 24 hour period.
	\end{solution}
	
	\columnbreak
	
	\hfil \includegraphics[width=0.35\textwidth]{mixed_semidiurnal_tide} \hfill
	
\end{multicols}

\AnswerBox{0.01\textheight}{}

\question[5]
\begin{multicols}{2}
	This figure shows the blood salinity of a marine worm relative to sea water. Is this worm mostly an osmoconformer or an osmoregulator? Tell how you know.
	
	\begin{solution}
		Osmoconformer. Blood salinity changes with external salinity.
	\end{solution}

	\columnbreak
	
	\hfil \noindent\includegraphics[width=0.35\textwidth]{osmoconformer_graph}\hfill
	
\end{multicols}

\AnswerBox{0.01\textheight}{}

\question[5]
Why is a stenohaline marine fish in a coastal plain estuary more likely to be located deeper in the water column the farther it swims up the estuary (towards the main river)?

\AnswerBox{0.1\textheight}{Deep estuaries are more likely to form a salt water wedge. The denser sea water will be deeper than the less dense freshwater.}

%\question[5]
%How do hypersaline patches form in New England marshes?

%\AnswerBox{\basespace}{Wrack or ice opens bare patches. Evaporation removes water, leaving salt behind, causing hypersaline conditions.}


\bonusquestion[1]
Write the name of your favorite band or musical artist.


\newpage

\question[5]
Explain bioturbation and how it increases oxygen availability in the sediments of a soft-bottom intertidal community.

\AnswerBox{0.1\textheight}{Bioturbation is the movement of organisms through the substrate. Their movement increases water flow through the sediment. The water contains oxygen so more water means more oxygen available to the organisms.}

%\question[8]
%Two days ago was a full moon. Tell whether that phase would have caused a neap tide or a spring tide, and then diagram and explain why.

%\AnswerBox{2\basespace}{The sun and moon are in a direct line with Earth, increasing the gravitational pull. This causes higher high tides and lower low tides, which are spring tides.}

%\question[8]
%Three days ago was a half moon. Tell whether that phase would have caused a neap tide or a spring tide, and then diagram and explain why.
%
%\AnswerBox{2\basespace}{The sun and moon are at right angles to each other, relative to Earth, negating the gravitational pull of each. This causes lower high tides and higher low tides, which are neap tides.}

\question[5]
Briefly explain why newly created patches larger than 1m tend to increase diversity in the rocky intertidal. (\textsc{Hint:} California and Bay mussels.)

\AnswerBox{0.1\textheight}{Large patches give other spp a chance to colonize before original species fully recolonize.}

%\question[5]
%Briefly describe how the recruitment limitation model explains diversity of fishes on coral reef.
%
%\AnswerBox{0.1\textheight}{Recruits are limited\dots}
%
\question[8]
A half moon is approaching. Tell whether that phase would have caused a neap tide or a spring tide, and then diagram and explain why.

\AnswerBox{0.3\textheight}{The sun and moon are at right angles to each other, relative to Earth, negating the gravitational pull of each. This causes lower high tides and higher low tides, which are neap tides.}

%\question[10]
%Explain how patch size affects whether a patch in the rocky intertidal will return to the original species or if a new species will be able to colonize and remain established. (\textsc{Hint:} California and Bay mussels.)
%
%\AnswerBox{2\basespace}{Small patches less than 1m more likely to return to orignal.
%}

\question[10]
Compare and contrast the competition and recruitment limitation models of reef fish community structure in terms of competition and post-recruitment processes.

\AnswerBox{0.25\textheight}{Competition has strong comp and affected by post-recruitment processes. Recruitment limitation has weak comp and not modified by post-recruitment.}


%\question[5]
%Describe the size distribution of POM in a bar-built estuary from the river input to the sand bar. Justify your answer.
%
%\AnswerBox{0.25\textheight}{Larger POM in sandy areas. Fine POM in mud areas. POM settles out at same size as inorganic sediment.}

%\newpage

%\question[5]

%Name the two groups of non-coral organisms that contribute to biogenic reef building.

%\AnswerBox{0.01\textheight}{Coralline red algae, calcareous green algae.}

\question[5]
Describe how patch dynamics at local scales in a rocky intertidal affects diversity at regional scales.


\AnswerBox{0.1\textheight}{Answer goes here.}

\question[10]
Describe how the Coriolis effect causes salinity to vary from side to side. Diagram this for an estuary that runs north to south in the \emph{southern} hemisphere.

\AnswerBox{0.3\textheight}{Large near freshwater input because sand grains fall out first, leaving larger grains so larger individuals required to ingest grains. In the middle, clay and silt fall out, so smaller deposit feeders. Behind the sandbar, wave action removes the smaller sediments, leaving sand so once again larger deposit feeders.}

\question[15]
Deposit feeders are small animals that ingest sediments to digest any organic material present between sediment grains. Describe how the relative size (larger, smaller, etc.) of the deposit feeders change in a bar-built estuary from the river to the sand bar at the ocean. Tell why the size of the deposit feeders changes. Deposit feeder size is proportional to sediment size. 

\AnswerBox{0.4\textheight}{Large near freshwater input because sand grains fall out first, leaving larger grains so larger individuals required to ingest grains. In the middle, clay and silt fall out, so smaller deposit feeders. Behind the sandbar, wave action removes the smaller sediments, leaving sand so once again larger deposit feeders.}

%\question[15]
%Meiofauna refers to the small animals that live between sediment grains. Describe how the relative size (larger, smaller, etc.) of the meiofauna will change in a bar-built estuary from the river to the sand bar at the ocean. Tell why the size of the meiofauna changes. % Tell how water movement (river flow or wave action) determines this distribution.
%
%\AnswerBox{0.7\textheight}{Large near freshwater input because sand grains fall out first, leaving larger spaces between the grains. In the middle, clay and silt fall out, so small spaces and thus smaller meiofauna. Behind , smaller with mud in middle of estuary, then large again behind the sand bar.}

\newpage

%\question[15]
%Compare and contrast the abiotic and biotic processes that determine zonation in a rocky intertidal compared to a salt marsh. Be sure to describe which processes (e.g., competition, wave action, etc.) is most important the upper and lower boundaries of each zone.
%
%\AnswerBox{0.5\textheight}{Larger POM in sandy areas. Fine POM in mud areas. POM settles out at same size as inorganic sediment.}


%\question[15]
%Name the four major zones found in a typical New England salt marsh, based on the characteristic plant species. Explain how physical and competitive forces interact to create these zones.

%\question
%	\begin{parts}
%	\part[5]
%	Illustrate the profile of a barrier reef from shoreline to open ocean. Label each zone.
%
%	\AnswerBox{0.3\textheight}{%
%	Fore-reef slope, reef crest, reef flat.%
%}
%
%
%	\part[10]
%	Describe how light, aerial exposure, and wave action interact to determine the types of hermatypic corals (relative size, branching, etc.) present in each zone.
%	
%		\begin{AnswerPage}{0.5\textheight}
%		Answer goes here.
%		\end{AnswerPage}
%	\end{parts}

\end{questions}

\end{document}